\section{Discussion}
\label{sec:discussion}
\subsection{Detection quality and transfer learning}
Among the runs evaluated over 10,260 validation frames, fine-tuned YOLOv8n has the highest reported AP$_{50}$ and AP$_{50:95}$, followed by the ImageNet-initialized Faster R-CNN variant. This observation is limited to the published runs. The Faster R-CNN variants differ in both backbone architecture and initialization, and the YOLO experiment uses a different training pipeline and sampling configuration. The results therefore do not isolate the contribution of pretraining, architecture or training data volume.

The fine-tuned YOLOv8n AP$_{50:95}$ of 0.564 exceeds the COCO baseline's 0.259 in the reported validation evaluation. Fine-tuning is a plausible contributor, but the baseline also lacks a matching pretrained van category. The difference cannot be attributed solely to weight adaptation. The compact RT-DETR-inspired run obtains high recall at confidence 0.25 but substantially lower precision and AP, indicating that detections at that operating point contain many false positives; the result warrants inspection by class, distance and confidence rather than a claim about official RT-DETR \cite{comparison}.

\subsection{Small vehicles and measurement geometry}
The published AP for small objects differs strongly across runs: 0.419 for fine-tuned YOLOv8n, 0.232 for scratch Faster R-CNN and 0.056 for its ImageNet/ResNet-18 variant. These are evaluator categories and do not automatically correspond to a particular polygon or camera distance. Zone-specific recall, object sizes and a consistent per-sequence evaluation would be needed to quantify the proposed far-zone strategy. Similarly, confidence thresholds and crop inference affect recall and false positives and must be selected on validation data.

Oblique road views remain geometrically usable if a sufficiently visible road region is approximately planar and metric correspondences are known. Near the horizon, however, perspective amplifies image-coordinate errors; low camera angles and occlusion can make the bottom-centre of a box an unreliable proxy for road contact. A Kalman filter can model some observation noise, but it cannot correct a wrong physical scale or systematic contact-point bias. Calibration error must be measured on independent road points, and a speed claim requires ground-truth distance/time or an external speed measurement.

\subsection{Limitations and next steps}
\label{sec:limitations}
The available detection results are validation reports, not held-out test results. Fast R-CNN and R-CNN were measured on a smaller subset. Checkpoint selection, step caps and different training budgets further limit causal comparison. The training figure mixes different model losses and validation procedures, and its series should not be read as a common-scale benchmark. The available artifacts inspected for this report do not establish HOTA, IDF1, MOTA or ID switches for the final system, nor MAE/RMSE against ground-truth speed. The supplied evaluation code supports such metrics when appropriate reference data are provided; implementation is not evidence of achieved performance.

Future work should (1) run each selected detector on the same untouched test sequences, with class and ignore policies fixed; (2) evaluate tracking against trajectory IDs using the same detections and an explicitly documented protocol; (3) validate each camera's road-plane calibration and measure speed against independent ground truth; (4) benchmark latency and memory on a named device with a fixed input and timing protocol; and (5) repeat experiments with multiple seeds where feasible. Camera configuration, actual training sample counts, code revision and environment should be archived with each run.

\section{Conclusion}
TrafficVision combines vehicle detection, interchangeable tracking, polygon zones and a calibrated road-plane speed pipeline. The inspected project artifacts include detector checkpoints and validation results: fine-tuned YOLOv8n reports AP$_{50}=0.750$, while the ImageNet-initialized Faster R-CNN variant reports 0.719 over the published 10,260 validation frames. These results demonstrate functioning detector experiments but do not establish final test performance or verified physical speed accuracy. A common held-out protocol and independent camera/speed measurements are the next steps for a defensible end-to-end assessment.
